\documentclass[11pt]{article}
% Official ACL style (github.com/acl-org/acl-style-files; commit in ACL_STYLE_VERSION.txt).
% "preprint" is non-anonymous with page numbers; use "review" for anonymous submission.
\usepackage[preprint]{acl}
\usepackage{times}
\usepackage{latexsym}
\usepackage[T1]{fontenc}
\usepackage[utf8]{inputenc}
\usepackage{microtype}
\usepackage{inconsolata}
\usepackage{amsmath,amssymb,booktabs,array,xcolor}
\newcommand{\F}{\mathcal{F}}
\newcommand{\D}{\mathcal{D}}
\newcommand{\ind}{\mathbf{1}}
\newcommand{\code}[1]{\texttt{#1}}
% Results that do not yet exist render visibly; never replace these by hand.
\newcommand{\pending}{\textcolor{red}{\textbf{[pending]}}}
% Prose that must be rewritten after the held-out results are known.
\newcommand{\afterresults}[1]{\textcolor{red}{[After results: #1]}}
\hypersetup{pdftitle={EchoTrace: Do Language Models Mistake Repeated Reports for Independent Corroboration?},pdfauthor={Ayush Dev}}
\title{EchoTrace: Do Language Models Mistake Repeated Reports for Independent Corroboration?}
\author{Ayush Dev\\
  Nepal Engineering College\\
  Bhaktapur, Nepal}
% Generated from frozen v1.1 test artifacts; do not edit manually.
\newcommand{\UrlFOne}{0.300}
\newcommand{\UrlMAE}{1.600}
\newcommand{\OneFOne}{0.507}
\newcommand{\OneMAE}{1.400}
\newcommand{\DomainFOne}{0.200}
\newcommand{\DomainMAE}{1.900}
\newcommand{\BowFOne}{0.507}
\newcommand{\BowMAE}{1.400}
\newcommand{\TfidfFOne}{0.467}
\newcommand{\TfidfMAE}{1.100}
\newcommand{\MinilmFOne}{0.507}
\newcommand{\MinilmMAE}{1.400}
\newcommand{\EchoFOne}{0.734}
\newcommand{\EchoMAE}{0.455}
\newcommand{\NoLinksFOne}{0.635}
\newcommand{\NoLinksMAE}{0.625}
\newcommand{\NoAttrFOne}{0.731}
\newcommand{\NoAttrMAE}{0.465}
\newcommand{\NoLexFOne}{0.734}
\newcommand{\NoLexMAE}{0.455}
\newcommand{\NoNumFOne}{0.792}
\newcommand{\NoNumMAE}{0.447}
\newcommand{\NoChronFOne}{0.734}
\newcommand{\NoChronMAE}{0.455}
\newcommand{\NoDomainFOne}{0.734}
\newcommand{\NoDomainMAE}{0.455}
\newcommand{\EchoFOneCI}{[0.704, 0.762]}
\newcommand{\EchoMAECI}{[0.378, 0.535]}
\newcommand{\TestCases}{600}
\newcommand{\TestWorlds}{60}
\newcommand{\DocProbErr}{0.082}
\newcommand{\DocLogErr}{1.015}
\newcommand{\SimTfidfProbErr}{0.056}
\newcommand{\SimTfidfLogErr}{0.698}
\newcommand{\SimMinilmProbErr}{0.085}
\newcommand{\SimMinilmLogErr}{0.888}
\newcommand{\SimEchoProbErr}{0.029}
\newcommand{\SimEchoLogErr}{0.357}
\newcommand{\EchoDocProbDelta}{-0.052}
\newcommand{\EchoDocProbDeltaCI}{[-0.063, -0.042]}
\newcommand{\PairedUrlFOne}{0.434}
\newcommand{\PairedUrlFOneCI}{[0.404, 0.462]}
\newcommand{\PairedUrlMAE}{-1.145}
\newcommand{\PairedUrlMAECI}{[-1.218, -1.067]}
\newcommand{\PairedTfidfFOne}{0.267}
\newcommand{\PairedTfidfFOneCI}{[0.237, 0.295]}
\newcommand{\PairedTfidfMAE}{-0.645}
\newcommand{\PairedTfidfMAECI}{[-0.718, -0.567]}
\newcommand{\ChronEdges}{2{,}250}
\newcommand{\ChronReversed}{577}

% Generated by scripts/export_llm_study_results.py; do not edit manually.
\newcommand{\StudyStatus}{\pending}
\newcommand{\QwenCoverage}{\pending}
\newcommand{\QwenStandardFcr}{\pending}
\newcommand{\QwenStandardFcrCI}{\pending}
\newcommand{\QwenStandardSens}{\pending}
\newcommand{\QwenStandardSensCI}{\pending}
\newcommand{\QwenStandardAcc}{\pending}
\newcommand{\QwenStandardUnder}{\pending}
\newcommand{\QwenStandardDetAbst}{\pending}
\newcommand{\QwenStandardCeil}{\pending}
\newcommand{\QwenStandardN}{\pending}
\newcommand{\QwenStandardAbst}{\pending}
\newcommand{\QwenStandardCountErr}{\pending}
\newcommand{\QwenReminderFcr}{\pending}
\newcommand{\QwenReminderFcrCI}{\pending}
\newcommand{\QwenReminderSens}{\pending}
\newcommand{\QwenReminderSensCI}{\pending}
\newcommand{\QwenReminderAcc}{\pending}
\newcommand{\QwenReminderUnder}{\pending}
\newcommand{\QwenReminderDetAbst}{\pending}
\newcommand{\QwenReminderCeil}{\pending}
\newcommand{\QwenReminderN}{\pending}
\newcommand{\QwenReminderAbst}{\pending}
\newcommand{\QwenReminderCountErr}{\pending}
\newcommand{\QwenEchographFcr}{\pending}
\newcommand{\QwenEchographFcrCI}{\pending}
\newcommand{\QwenEchographSens}{\pending}
\newcommand{\QwenEchographSensCI}{\pending}
\newcommand{\QwenEchographAcc}{\pending}
\newcommand{\QwenEchographUnder}{\pending}
\newcommand{\QwenEchographDetAbst}{\pending}
\newcommand{\QwenEchographCeil}{\pending}
\newcommand{\QwenEchographN}{\pending}
\newcommand{\QwenEchographAbst}{\pending}
\newcommand{\QwenEchographCountErr}{\pending}
\newcommand{\QwenOracleFcr}{\pending}
\newcommand{\QwenOracleFcrCI}{\pending}
\newcommand{\QwenOracleSens}{\pending}
\newcommand{\QwenOracleSensCI}{\pending}
\newcommand{\QwenOracleAcc}{\pending}
\newcommand{\QwenOracleUnder}{\pending}
\newcommand{\QwenOracleDetAbst}{\pending}
\newcommand{\QwenOracleCeil}{\pending}
\newcommand{\QwenOracleN}{\pending}
\newcommand{\QwenOracleAbst}{\pending}
\newcommand{\QwenOracleCountErr}{\pending}
\newcommand{\QwenEchoStdFcr}{\pending}
\newcommand{\QwenEchoStdFcrCI}{\pending}
\newcommand{\QwenEchoStdAcc}{\pending}
\newcommand{\QwenEchoStdAccCI}{\pending}
\newcommand{\QwenEchoStdUnder}{\pending}
\newcommand{\QwenEchoStdUnderCI}{\pending}
\newcommand{\QwenRemStdFcr}{\pending}
\newcommand{\QwenRemStdFcrCI}{\pending}
\newcommand{\QwenRemStdAcc}{\pending}
\newcommand{\QwenRemStdAccCI}{\pending}
\newcommand{\QwenRemStdUnder}{\pending}
\newcommand{\QwenRemStdUnderCI}{\pending}
\newcommand{\QwenEchoRemFcr}{\pending}
\newcommand{\QwenEchoRemFcrCI}{\pending}
\newcommand{\QwenEchoRemAcc}{\pending}
\newcommand{\QwenEchoRemAccCI}{\pending}
\newcommand{\QwenEchoRemUnder}{\pending}
\newcommand{\QwenEchoRemUnderCI}{\pending}
\newcommand{\QwenOracleEchoFcr}{\pending}
\newcommand{\QwenOracleEchoFcrCI}{\pending}
\newcommand{\QwenOracleEchoAcc}{\pending}
\newcommand{\QwenOracleEchoAccCI}{\pending}
\newcommand{\QwenOracleEchoUnder}{\pending}
\newcommand{\QwenOracleEchoUnderCI}{\pending}
\newcommand{\QwenHolmP}{\pending}
\newcommand{\GptossCoverage}{\pending}
\newcommand{\GptossStandardFcr}{\pending}
\newcommand{\GptossStandardFcrCI}{\pending}
\newcommand{\GptossStandardSens}{\pending}
\newcommand{\GptossStandardSensCI}{\pending}
\newcommand{\GptossStandardAcc}{\pending}
\newcommand{\GptossStandardUnder}{\pending}
\newcommand{\GptossStandardDetAbst}{\pending}
\newcommand{\GptossStandardCeil}{\pending}
\newcommand{\GptossStandardN}{\pending}
\newcommand{\GptossStandardAbst}{\pending}
\newcommand{\GptossStandardCountErr}{\pending}
\newcommand{\GptossReminderFcr}{\pending}
\newcommand{\GptossReminderFcrCI}{\pending}
\newcommand{\GptossReminderSens}{\pending}
\newcommand{\GptossReminderSensCI}{\pending}
\newcommand{\GptossReminderAcc}{\pending}
\newcommand{\GptossReminderUnder}{\pending}
\newcommand{\GptossReminderDetAbst}{\pending}
\newcommand{\GptossReminderCeil}{\pending}
\newcommand{\GptossReminderN}{\pending}
\newcommand{\GptossReminderAbst}{\pending}
\newcommand{\GptossReminderCountErr}{\pending}
\newcommand{\GptossEchographFcr}{\pending}
\newcommand{\GptossEchographFcrCI}{\pending}
\newcommand{\GptossEchographSens}{\pending}
\newcommand{\GptossEchographSensCI}{\pending}
\newcommand{\GptossEchographAcc}{\pending}
\newcommand{\GptossEchographUnder}{\pending}
\newcommand{\GptossEchographDetAbst}{\pending}
\newcommand{\GptossEchographCeil}{\pending}
\newcommand{\GptossEchographN}{\pending}
\newcommand{\GptossEchographAbst}{\pending}
\newcommand{\GptossEchographCountErr}{\pending}
\newcommand{\GptossOracleFcr}{\pending}
\newcommand{\GptossOracleFcrCI}{\pending}
\newcommand{\GptossOracleSens}{\pending}
\newcommand{\GptossOracleSensCI}{\pending}
\newcommand{\GptossOracleAcc}{\pending}
\newcommand{\GptossOracleUnder}{\pending}
\newcommand{\GptossOracleDetAbst}{\pending}
\newcommand{\GptossOracleCeil}{\pending}
\newcommand{\GptossOracleN}{\pending}
\newcommand{\GptossOracleAbst}{\pending}
\newcommand{\GptossOracleCountErr}{\pending}
\newcommand{\GptossEchoStdFcr}{\pending}
\newcommand{\GptossEchoStdFcrCI}{\pending}
\newcommand{\GptossEchoStdAcc}{\pending}
\newcommand{\GptossEchoStdAccCI}{\pending}
\newcommand{\GptossEchoStdUnder}{\pending}
\newcommand{\GptossEchoStdUnderCI}{\pending}
\newcommand{\GptossRemStdFcr}{\pending}
\newcommand{\GptossRemStdFcrCI}{\pending}
\newcommand{\GptossRemStdAcc}{\pending}
\newcommand{\GptossRemStdAccCI}{\pending}
\newcommand{\GptossRemStdUnder}{\pending}
\newcommand{\GptossRemStdUnderCI}{\pending}
\newcommand{\GptossEchoRemFcr}{\pending}
\newcommand{\GptossEchoRemFcrCI}{\pending}
\newcommand{\GptossEchoRemAcc}{\pending}
\newcommand{\GptossEchoRemAccCI}{\pending}
\newcommand{\GptossEchoRemUnder}{\pending}
\newcommand{\GptossEchoRemUnderCI}{\pending}
\newcommand{\GptossOracleEchoFcr}{\pending}
\newcommand{\GptossOracleEchoFcrCI}{\pending}
\newcommand{\GptossOracleEchoAcc}{\pending}
\newcommand{\GptossOracleEchoAccCI}{\pending}
\newcommand{\GptossOracleEchoUnder}{\pending}
\newcommand{\GptossOracleEchoUnderCI}{\pending}
\newcommand{\GptossHolmP}{\pending}
\newcommand{\GeminiCoverage}{\pending}
\newcommand{\GeminiStandardFcr}{\pending}
\newcommand{\GeminiStandardFcrCI}{\pending}
\newcommand{\GeminiStandardSens}{\pending}
\newcommand{\GeminiStandardSensCI}{\pending}
\newcommand{\GeminiStandardAcc}{\pending}
\newcommand{\GeminiStandardUnder}{\pending}
\newcommand{\GeminiStandardDetAbst}{\pending}
\newcommand{\GeminiStandardCeil}{\pending}
\newcommand{\GeminiStandardN}{\pending}
\newcommand{\GeminiStandardAbst}{\pending}
\newcommand{\GeminiStandardCountErr}{\pending}
\newcommand{\GeminiReminderFcr}{\pending}
\newcommand{\GeminiReminderFcrCI}{\pending}
\newcommand{\GeminiReminderSens}{\pending}
\newcommand{\GeminiReminderSensCI}{\pending}
\newcommand{\GeminiReminderAcc}{\pending}
\newcommand{\GeminiReminderUnder}{\pending}
\newcommand{\GeminiReminderDetAbst}{\pending}
\newcommand{\GeminiReminderCeil}{\pending}
\newcommand{\GeminiReminderN}{\pending}
\newcommand{\GeminiReminderAbst}{\pending}
\newcommand{\GeminiReminderCountErr}{\pending}
\newcommand{\GeminiEchographFcr}{\pending}
\newcommand{\GeminiEchographFcrCI}{\pending}
\newcommand{\GeminiEchographSens}{\pending}
\newcommand{\GeminiEchographSensCI}{\pending}
\newcommand{\GeminiEchographAcc}{\pending}
\newcommand{\GeminiEchographUnder}{\pending}
\newcommand{\GeminiEchographDetAbst}{\pending}
\newcommand{\GeminiEchographCeil}{\pending}
\newcommand{\GeminiEchographN}{\pending}
\newcommand{\GeminiEchographAbst}{\pending}
\newcommand{\GeminiEchographCountErr}{\pending}
\newcommand{\GeminiOracleFcr}{\pending}
\newcommand{\GeminiOracleFcrCI}{\pending}
\newcommand{\GeminiOracleSens}{\pending}
\newcommand{\GeminiOracleSensCI}{\pending}
\newcommand{\GeminiOracleAcc}{\pending}
\newcommand{\GeminiOracleUnder}{\pending}
\newcommand{\GeminiOracleDetAbst}{\pending}
\newcommand{\GeminiOracleCeil}{\pending}
\newcommand{\GeminiOracleN}{\pending}
\newcommand{\GeminiOracleAbst}{\pending}
\newcommand{\GeminiOracleCountErr}{\pending}
\newcommand{\GeminiEchoStdFcr}{\pending}
\newcommand{\GeminiEchoStdFcrCI}{\pending}
\newcommand{\GeminiEchoStdAcc}{\pending}
\newcommand{\GeminiEchoStdAccCI}{\pending}
\newcommand{\GeminiEchoStdUnder}{\pending}
\newcommand{\GeminiEchoStdUnderCI}{\pending}
\newcommand{\GeminiRemStdFcr}{\pending}
\newcommand{\GeminiRemStdFcrCI}{\pending}
\newcommand{\GeminiRemStdAcc}{\pending}
\newcommand{\GeminiRemStdAccCI}{\pending}
\newcommand{\GeminiRemStdUnder}{\pending}
\newcommand{\GeminiRemStdUnderCI}{\pending}
\newcommand{\GeminiEchoRemFcr}{\pending}
\newcommand{\GeminiEchoRemFcrCI}{\pending}
\newcommand{\GeminiEchoRemAcc}{\pending}
\newcommand{\GeminiEchoRemAccCI}{\pending}
\newcommand{\GeminiEchoRemUnder}{\pending}
\newcommand{\GeminiEchoRemUnderCI}{\pending}
\newcommand{\GeminiOracleEchoFcr}{\pending}
\newcommand{\GeminiOracleEchoFcrCI}{\pending}
\newcommand{\GeminiOracleEchoAcc}{\pending}
\newcommand{\GeminiOracleEchoAccCI}{\pending}
\newcommand{\GeminiOracleEchoUnder}{\pending}
\newcommand{\GeminiOracleEchoUnderCI}{\pending}
\newcommand{\GeminiHolmP}{\pending}
\newcommand{\DevStudyStatus}{complete}
\newcommand{\DevQwenCoverage}{0.708}
\newcommand{\DevQwenStandardFcr}{0.037}
\newcommand{\DevQwenStandardFcrCI}{[0.000, 0.121]}
\newcommand{\DevQwenStandardSens}{0.006}
\newcommand{\DevQwenStandardSensCI}{[0.000, 0.013]}
\newcommand{\DevQwenStandardAcc}{0.902}
\newcommand{\DevQwenStandardUnder}{0.000}
\newcommand{\DevQwenStandardDetAbst}{0.000}
\newcommand{\DevQwenStandardCeil}{0.896}
\newcommand{\DevQwenStandardN}{51}
\newcommand{\DevQwenStandardAbst}{0.625}
\newcommand{\DevQwenStandardCountErr}{0.023}
\newcommand{\DevQwenReminderFcr}{0.037}
\newcommand{\DevQwenReminderFcrCI}{[0.000, 0.121]}
\newcommand{\DevQwenReminderSens}{0.006}
\newcommand{\DevQwenReminderSensCI}{[0.000, 0.013]}
\newcommand{\DevQwenReminderAcc}{0.882}
\newcommand{\DevQwenReminderUnder}{0.062}
\newcommand{\DevQwenReminderDetAbst}{0.023}
\newcommand{\DevQwenReminderCeil}{0.438}
\newcommand{\DevQwenReminderN}{51}
\newcommand{\DevQwenReminderAbst}{0.750}
\newcommand{\DevQwenReminderCountErr}{0.095}
\newcommand{\DevQwenEchographFcr}{0.000}
\newcommand{\DevQwenEchographFcrCI}{[0.000, 0.000]}
\newcommand{\DevQwenEchographSens}{0.000}
\newcommand{\DevQwenEchographSensCI}{[0.000, 0.000]}
\newcommand{\DevQwenEchographAcc}{0.863}
\newcommand{\DevQwenEchographUnder}{0.000}
\newcommand{\DevQwenEchographDetAbst}{0.047}
\newcommand{\DevQwenEchographCeil}{0.542}
\newcommand{\DevQwenEchographN}{51}
\newcommand{\DevQwenEchographAbst}{0.375}
\newcommand{\DevQwenEchographCountErr}{0.000}
\newcommand{\DevQwenOracleFcr}{0.000}
\newcommand{\DevQwenOracleFcrCI}{[0.000, 0.000]}
\newcommand{\DevQwenOracleSens}{0.000}
\newcommand{\DevQwenOracleSensCI}{[0.000, 0.000]}
\newcommand{\DevQwenOracleAcc}{0.865}
\newcommand{\DevQwenOracleUnder}{0.059}
\newcommand{\DevQwenOracleDetAbst}{0.023}
\newcommand{\DevQwenOracleCeil}{0.521}
\newcommand{\DevQwenOracleN}{52}
\newcommand{\DevQwenOracleAbst}{0.750}
\newcommand{\DevQwenOracleCountErr}{0.070}
\newcommand{\DevQwenEchoStdFcr}{-0.037}
\newcommand{\DevQwenEchoStdFcrCI}{[-0.121, 0.000]}
\newcommand{\DevQwenEchoStdAcc}{-0.039}
\newcommand{\DevQwenEchoStdAccCI}{[-0.093, 0.000]}
\newcommand{\DevQwenEchoStdUnder}{0.000}
\newcommand{\DevQwenEchoStdUnderCI}{[0.000, 0.000]}
\newcommand{\DevQwenRemStdFcr}{0.000}
\newcommand{\DevQwenRemStdFcrCI}{[0.000, 0.000]}
\newcommand{\DevQwenRemStdAcc}{-0.020}
\newcommand{\DevQwenRemStdAccCI}{[-0.062, 0.000]}
\newcommand{\DevQwenRemStdUnder}{0.062}
\newcommand{\DevQwenRemStdUnderCI}{[0.000, 0.200]}
\newcommand{\DevQwenEchoRemFcr}{-0.037}
\newcommand{\DevQwenEchoRemFcrCI}{[-0.125, 0.000]}
\newcommand{\DevQwenEchoRemAcc}{-0.020}
\newcommand{\DevQwenEchoRemAccCI}{[-0.062, 0.000]}
\newcommand{\DevQwenEchoRemUnder}{-0.062}
\newcommand{\DevQwenEchoRemUnderCI}{[-0.200, 0.000]}
\newcommand{\DevQwenOracleEchoFcr}{0.000}
\newcommand{\DevQwenOracleEchoFcrCI}{[0.000, 0.000]}
\newcommand{\DevQwenOracleEchoAcc}{0.000}
\newcommand{\DevQwenOracleEchoAccCI}{[0.000, 0.000]}
\newcommand{\DevQwenOracleEchoUnder}{0.062}
\newcommand{\DevQwenOracleEchoUnderCI}{[0.000, 0.200]}
\newcommand{\DevQwenHolmP}{0.714}
\newcommand{\DevGptossCoverage}{0.764}
\newcommand{\DevGptossStandardFcr}{0.467}
\newcommand{\DevGptossStandardFcrCI}{[0.273, 0.667]}
\newcommand{\DevGptossStandardSens}{0.004}
\newcommand{\DevGptossStandardSensCI}{[-0.005, 0.013]}
\newcommand{\DevGptossStandardAcc}{0.877}
\newcommand{\DevGptossStandardUnder}{0.000}
\newcommand{\DevGptossStandardDetAbst}{0.021}
\newcommand{\DevGptossStandardCeil}{0.312}
\newcommand{\DevGptossStandardN}{57}
\newcommand{\DevGptossStandardAbst}{0.222}
\newcommand{\DevGptossStandardCountErr}{0.532}
\newcommand{\DevGptossReminderFcr}{0.000}
\newcommand{\DevGptossReminderFcrCI}{[0.000, 0.000]}
\newcommand{\DevGptossReminderSens}{-0.034}
\newcommand{\DevGptossReminderSensCI}{[-0.092, 0.020]}
\newcommand{\DevGptossReminderAcc}{0.776}
\newcommand{\DevGptossReminderUnder}{0.526}
\newcommand{\DevGptossReminderDetAbst}{0.265}
\newcommand{\DevGptossReminderCeil}{0.021}
\newcommand{\DevGptossReminderN}{58}
\newcommand{\DevGptossReminderAbst}{0.667}
\newcommand{\DevGptossReminderCountErr}{1.194}
\newcommand{\DevGptossEchographFcr}{0.000}
\newcommand{\DevGptossEchographFcrCI}{[0.000, 0.000]}
\newcommand{\DevGptossEchographSens}{-0.048}
\newcommand{\DevGptossEchographSensCI}{[-0.104, 0.012]}
\newcommand{\DevGptossEchographAcc}{0.804}
\newcommand{\DevGptossEchographUnder}{0.667}
\newcommand{\DevGptossEchographDetAbst}{0.170}
\newcommand{\DevGptossEchographCeil}{0.104}
\newcommand{\DevGptossEchographN}{56}
\newcommand{\DevGptossEchographAbst}{0.333}
\newcommand{\DevGptossEchographCountErr}{1.385}
\newcommand{\DevGptossOracleFcr}{0.000}
\newcommand{\DevGptossOracleFcrCI}{[0.000, 0.000]}
\newcommand{\DevGptossOracleSens}{-0.003}
\newcommand{\DevGptossOracleSensCI}{[-0.038, 0.033]}
\newcommand{\DevGptossOracleAcc}{0.818}
\newcommand{\DevGptossOracleUnder}{0.722}
\newcommand{\DevGptossOracleDetAbst}{0.109}
\newcommand{\DevGptossOracleCeil}{0.043}
\newcommand{\DevGptossOracleN}{55}
\newcommand{\DevGptossOracleAbst}{0.556}
\newcommand{\DevGptossOracleCountErr}{1.439}
\newcommand{\DevGptossEchoStdFcr}{-0.448}
\newcommand{\DevGptossEchoStdFcrCI}{[-0.639, -0.259]}
\newcommand{\DevGptossEchoStdAcc}{-0.071}
\newcommand{\DevGptossEchoStdAccCI}{[-0.152, 0.000]}
\newcommand{\DevGptossEchoStdUnder}{0.667}
\newcommand{\DevGptossEchoStdUnderCI}{[0.375, 0.929]}
\newcommand{\DevGptossRemStdFcr}{-0.467}
\newcommand{\DevGptossRemStdFcrCI}{[-0.667, -0.273]}
\newcommand{\DevGptossRemStdAcc}{-0.088}
\newcommand{\DevGptossRemStdAccCI}{[-0.167, -0.020]}
\newcommand{\DevGptossRemStdUnder}{0.500}
\newcommand{\DevGptossRemStdUnderCI}{[0.222, 0.778]}
\newcommand{\DevGptossEchoRemFcr}{0.000}
\newcommand{\DevGptossEchoRemFcrCI}{[0.000, 0.000]}
\newcommand{\DevGptossEchoRemAcc}{0.000}
\newcommand{\DevGptossEchoRemAccCI}{[-0.062, 0.048]}
\newcommand{\DevGptossEchoRemUnder}{0.167}
\newcommand{\DevGptossEchoRemUnderCI}{[-0.167, 0.550]}
\newcommand{\DevGptossOracleEchoFcr}{0.000}
\newcommand{\DevGptossOracleEchoFcrCI}{[0.000, 0.000]}
\newcommand{\DevGptossOracleEchoAcc}{0.019}
\newcommand{\DevGptossOracleEchoAccCI}{[-0.041, 0.093]}
\newcommand{\DevGptossOracleEchoUnder}{0.056}
\newcommand{\DevGptossOracleEchoUnderCI}{[0.000, 0.182]}
\newcommand{\DevGptossHolmP}{0.000}
\newcommand{\DevGeminiCoverage}{1.000}
\newcommand{\DevGeminiStandardFcr}{0.056}
\newcommand{\DevGeminiStandardFcrCI}{[0.000, 0.139]}
\newcommand{\DevGeminiStandardSens}{0.000}
\newcommand{\DevGeminiStandardSensCI}{[0.000, 0.000]}
\newcommand{\DevGeminiStandardAcc}{0.708}
\newcommand{\DevGeminiStandardUnder}{0.000}
\newcommand{\DevGeminiStandardDetAbst}{0.033}
\newcommand{\DevGeminiStandardCeil}{0.979}
\newcommand{\DevGeminiStandardN}{72}
\newcommand{\DevGeminiStandardAbst}{0.833}
\newcommand{\DevGeminiStandardCountErr}{0.052}
\newcommand{\DevGeminiReminderFcr}{0.000}
\newcommand{\DevGeminiReminderFcrCI}{[0.000, 0.000]}
\newcommand{\DevGeminiReminderSens}{0.000}
\newcommand{\DevGeminiReminderSensCI}{[0.000, 0.000]}
\newcommand{\DevGeminiReminderAcc}{0.639}
\newcommand{\DevGeminiReminderUnder}{0.000}
\newcommand{\DevGeminiReminderDetAbst}{0.250}
\newcommand{\DevGeminiReminderCeil}{0.896}
\newcommand{\DevGeminiReminderN}{72}
\newcommand{\DevGeminiReminderAbst}{1.000}
\newcommand{\DevGeminiReminderCountErr}{0.000}
\newcommand{\DevGeminiEchographFcr}{0.000}
\newcommand{\DevGeminiEchographFcrCI}{[0.000, 0.000]}
\newcommand{\DevGeminiEchographSens}{0.000}
\newcommand{\DevGeminiEchographSensCI}{[0.000, 0.000]}
\newcommand{\DevGeminiEchographAcc}{0.722}
\newcommand{\DevGeminiEchographUnder}{0.000}
\newcommand{\DevGeminiEchographDetAbst}{0.217}
\newcommand{\DevGeminiEchographCeil}{0.938}
\newcommand{\DevGeminiEchographN}{72}
\newcommand{\DevGeminiEchographAbst}{1.000}
\newcommand{\DevGeminiEchographCountErr}{0.000}
\newcommand{\DevGeminiOracleFcr}{0.000}
\newcommand{\DevGeminiOracleFcrCI}{[0.000, 0.000]}
\newcommand{\DevGeminiOracleSens}{0.006}
\newcommand{\DevGeminiOracleSensCI}{[0.000, 0.013]}
\newcommand{\DevGeminiOracleAcc}{0.694}
\newcommand{\DevGeminiOracleUnder}{0.000}
\newcommand{\DevGeminiOracleDetAbst}{0.267}
\newcommand{\DevGeminiOracleCeil}{0.896}
\newcommand{\DevGeminiOracleN}{72}
\newcommand{\DevGeminiOracleAbst}{1.000}
\newcommand{\DevGeminiOracleCountErr}{0.000}
\newcommand{\DevGeminiEchoStdFcr}{-0.056}
\newcommand{\DevGeminiEchoStdFcrCI}{[-0.139, 0.000]}
\newcommand{\DevGeminiEchoStdAcc}{0.014}
\newcommand{\DevGeminiEchoStdAccCI}{[-0.069, 0.125]}
\newcommand{\DevGeminiEchoStdUnder}{0.000}
\newcommand{\DevGeminiEchoStdUnderCI}{[0.000, 0.000]}
\newcommand{\DevGeminiRemStdFcr}{-0.056}
\newcommand{\DevGeminiRemStdFcrCI}{[-0.139, 0.000]}
\newcommand{\DevGeminiRemStdAcc}{-0.069}
\newcommand{\DevGeminiRemStdAccCI}{[-0.111, -0.028]}
\newcommand{\DevGeminiRemStdUnder}{0.000}
\newcommand{\DevGeminiRemStdUnderCI}{[0.000, 0.000]}
\newcommand{\DevGeminiEchoRemFcr}{0.000}
\newcommand{\DevGeminiEchoRemFcrCI}{[0.000, 0.000]}
\newcommand{\DevGeminiEchoRemAcc}{0.083}
\newcommand{\DevGeminiEchoRemAccCI}{[0.000, 0.194]}
\newcommand{\DevGeminiEchoRemUnder}{0.000}
\newcommand{\DevGeminiEchoRemUnderCI}{[0.000, 0.000]}
\newcommand{\DevGeminiOracleEchoFcr}{0.000}
\newcommand{\DevGeminiOracleEchoFcrCI}{[0.000, 0.000]}
\newcommand{\DevGeminiOracleEchoAcc}{-0.028}
\newcommand{\DevGeminiOracleEchoAccCI}{[-0.097, 0.028]}
\newcommand{\DevGeminiOracleEchoUnder}{0.000}
\newcommand{\DevGeminiOracleEchoUnderCI}{[0.000, 0.000]}
\newcommand{\DevGeminiHolmP}{0.448}
\newcommand{\UrlFOneVal}{0.333}
\newcommand{\UrlMAEVal}{1.444}
\newcommand{\OneFOneVal}{0.452}
\newcommand{\OneMAEVal}{1.556}
\newcommand{\DomainFOneVal}{0.222}
\newcommand{\DomainMAEVal}{1.778}
\newcommand{\BowFOneVal}{0.258}
\newcommand{\BowMAEVal}{1.370}
\newcommand{\TfidfFOneVal}{0.369}
\newcommand{\TfidfMAEVal}{1.309}
\newcommand{\MinilmFOneVal}{0.450}
\newcommand{\MinilmMAEVal}{1.552}
\newcommand{\EchoFOneVal}{0.914}
\newcommand{\EchoMAEVal}{0.085}
\newcommand{\NoLinksFOneVal}{0.698}
\newcommand{\NoLinksMAEVal}{0.520}
\newcommand{\EchoSharedFOneVal}{0.325}
\newcommand{\EchoFOneValCI}{[0.902, 0.926]}
\newcommand{\UrlFOneTest}{0.333}
\newcommand{\UrlMAETest}{1.444}
\newcommand{\OneFOneTest}{0.452}
\newcommand{\OneMAETest}{1.556}
\newcommand{\DomainFOneTest}{0.222}
\newcommand{\DomainMAETest}{1.778}
\newcommand{\BowFOneTest}{0.276}
\newcommand{\BowMAETest}{1.377}
\newcommand{\TfidfFOneTest}{0.376}
\newcommand{\TfidfMAETest}{1.315}
\newcommand{\MinilmFOneTest}{0.448}
\newcommand{\MinilmMAETest}{1.546}
\newcommand{\EchoFOneTest}{0.920}
\newcommand{\EchoMAETest}{0.086}
\newcommand{\NoLinksFOneTest}{0.697}
\newcommand{\NoLinksMAETest}{0.509}
\newcommand{\EchoSharedFOneTest}{0.369}
\newcommand{\EchoFOneTestCI}{[0.907, 0.933]}
\newcommand{\ProbeAccFour}{0.182}


\begin{document}
\maketitle

\begin{abstract}
When several documents support a claim, a research assistant may treat them as independent
confirmation even if they all derive from one observation. We study this \emph{false
corroboration} in language models. We release EchoBench-Synthetic v1.4, a chronology-valid,
gold-by-construction benchmark of 3,960 cases from 360 world-disjoint micro-worlds, 36 topics, and
six prose styles, two of which appear only in held-out worlds. In an exploratory study on 12
validation worlds, three freely accessible models (Qwen3.8-27B, GPT-OSS-120B, and Gemini 3.5
Flash-Lite) answer under four conditions: no guidance, a dependence reminder, source families
proposed by EchoGraph, an interpretable CPU detector, and gold families. Models must report how
many independent sources the evidence contains, which lets us score overcounting automatically.
Without guidance, GPT-OSS overcounts dependent sources in \DevGptossStandardFcr\ of cases (95\%
CI \DevGptossStandardFcrCI), usually by reporting the number of documents; the other two models
rarely do. Guidance removes overcounting but over-corrects: GPT-OSS then undercounts genuinely
independent sources in \DevGptossEchographUnder\ of cases with EchoGraph families and
\DevGptossOracleUnder\ even with gold families, and Gemini abstains on many determinate cases. No
model's confidence in its verdict depends on how many independent sources support it. These are
development findings, not confirmatory tests; we freeze the held-out protocol that would test
them. EchoGraph itself reaches family-pair F1 \EchoFOneTest\ in a single held-out run after the
freeze but remains weak on shared datasets. All experiments run on CPU and free API tiers; no human validation is claimed.
\end{abstract}

\section{Introduction}
A research assistant may cite four pages and describe a finding as repeatedly confirmed. If three
pages summarize the first, the answer contains four documents but one informational origin. We call
this \emph{false corroboration}: apparent support increases because dependent reports are counted
as independent observations. Citation checks do not expose it, because every page genuinely states
the claim.

Source dependence has long mattered in data integration \citep{dong2009dependence}, and citation
networks can amplify claims through repeated reference \citep{greenberg2009citation}. Current
evaluations of attributed and deep-research systems measure correctness, citation support,
evidential force, and synthesis quality \citep{gao2023alce,du2025deepresearch,qian2026force}. They
do not directly measure whether a system knows how many \emph{distinct} observations its sources
represent.
Recent formal and empirical work shows that aggregating multiplied agent reports can destroy
calibration when their shared ancestry is ignored \citep{bara2026sybil}. Whether language models make this error when provenance must be read from the
documents themselves, and whether simple provenance metadata prevents it, has not been tested
directly.

We test it in a controlled setting where provenance is known by construction. Each case gives a
model the same four claim-relevant documents, varying only whether they share an origin. We ask
three questions:
\begin{enumerate}
  \item \textbf{False corroboration.} Without guidance, do models report more independent sources
  than exist when documents share an origin?
  \item \textbf{Confidence sensitivity.} Holding the claim and verdict fixed, is a model more
  confident when evidence comes from four independent sources than from one?
  \item \textbf{Mitigation.} Do source families proposed by an interpretable detector reduce false
  corroboration, beyond a plain reminder, without lowering verdict accuracy or causing models to
  undercount independent evidence?
\end{enumerate}

Our contributions are:
\begin{itemize}
  \item a zero-cost exploratory evaluation of three model families across four guidance
  conditions, with automatically scored source counts and whole-world bootstrap intervals,
  showing that guidance can trade overcounting for undercounting;
  \item a frozen held-out confirmatory protocol, amended from these development observations with
  an over-correction guard, that others can run under the same free quotas;
  \item EchoBench-Synthetic v1.4, with chronology-valid lineage, observable provenance cues,
  36 topics, and held-out prose styles; and
  \item EchoGraph, a CPU-only logistic source-family detector, evaluated against document,
  domain, lexical, TF--IDF, and sentence-embedding baselines, together with an account of a
  construction defect in our earlier release and its consequences.
\end{itemize}
The study measures how models use \emph{visible} provenance in synthetic documents. It does not
establish open-web provenance recovery, factual truth, or human-validated answer quality.

\section{Related Work}
\paragraph{Attribution and research-agent evaluation.}
ALCE evaluates correctness and citation quality in attributed answers \citep{gao2023alce}.
Deep-research benchmarks evaluate long reports and cross-source reasoning
\citep{du2025deepresearch,xie2026deepweb}, evidence-force evaluation separates relevance from
warrant \citep{qian2026force}, and audits show that displayed citations need not provide verified
attribution \citep{onweller2026cited}. Source choice can also change answers in multi-source medical
RAG \citep{li2026transplantqa}. EchoTrace measures a different property: relationships
\emph{between} supporting sources.

\paragraph{Dependence and diversity.}
Copy-detection models for data integration avoid double-counting repeated values
\citep{dong2009dependence}. MMR reduces redundancy by trading relevance against novelty
\citep{goldstein1998mmr}, but textual diversity is not evidential independence: independent studies
may use similar language and derivatives may paraphrase. LineageRAG builds evidence lineages grounded in
verbatim source spans \citep{zheng2026lineagerag}; EchoTrace instead infers which documents share
an origin.

\paragraph{Report multiplicity.}
\citet{bara2026sybil} formalizes when an additional report adds no evidence, and shows with
controlled LLM-agent extraction on synthetic documents that naive aggregation loses calibration as
reports of one evidence root multiply. It also finds that report similarity, not true ancestry,
drives a similarity-based deduplicator, which parallels our same-wording condition. That work
studies aggregation over agent reports. EchoTrace is complementary and narrower: it asks whether a
single model, reading the documents themselves, counts dependent reports as independent sources,
and whether provenance guidance changes that count in either direction.

\section{Task and Measures}
For a focal claim $c$, let $\D_c=\{d_1,\ldots,d_n\}$ be its relevant documents. In a determinate
case, a latent mapping $z(d)$ assigns each document to one source family, giving
\begin{equation}
  \F_c=\{F_k\}_{k=1}^{K},\quad F_k=\{d\in\D_c:z(d)=k\},
\end{equation}
where $K$ counts informational origins, not URLs or publishers. In an indeterminate case the
correct response is to abstain from a count.

\paragraph{Language-model measures.}
A model returns a verdict (supported, contradicted, insufficient), verdict probabilities, citations,
and $\widehat K$, the number of independent sources bearing on the claim, or \code{null} if this
cannot be determined. The \emph{false-corroboration rate} (FCR) is the fraction of responses on
dependent cases ($K<4$) with $\widehat K>K$; an abstention is not an overcount. \emph{Confidence sensitivity} is the within-world
difference in the probability assigned to the gold verdict between the four-independent and
single-origin versions of the same claim; a model that weighs independence should give a positive
value. We also report undercounting on independent cases ($\widehat K<K$), abstention on
determinate and on indeterminate cases, verdict accuracy, and the share of gold-verdict
probabilities at a 0.99 ceiling, which shows whether confidence could rise at all.

\paragraph{Detector measures.}
For detectors returning a partition $\widehat\F_c$, we report root-count mean absolute error
$\frac1N\sum_i|\widehat K_i-K_i|$ and exclusive-family pair F1 over same-family document pairs.

Source-family dependence is not truth or reliability. Two independent studies may share a bias,
and a derivative can accurately summarize strong evidence.

\section{EchoBench-Synthetic v1.4}
\paragraph{Construction.}
EchoBench contains 360 fictional micro-worlds, each with one quantitative claim and one verdict,
crossed with 11 provenance conditions (Table~\ref{tab:conditions}). Every case has 12 documents:
four claim-relevant documents and eight distractors. Whole worlds are assigned to train,
validation, and test (216/72/72), with verdicts balanced within each split. Documents use reserved
\code{.example} domains and opaque identifiers. Public files contain only queries and documents;
evaluator files contain relevance, verdicts, lineage, families, conditions, and splits.

\begin{table}[t]
\caption{Provenance conditions. $K$ is the gold number of source families among the four
relevant documents; ``--'' denotes an abstention target. $\star$ marks conditions used in the
language-model study.}
\label{tab:conditions}
\centering\footnotesize
\setlength{\tabcolsep}{3pt}
\begin{tabular}{@{}lc>{\raggedright\arraybackslash}p{2.9cm}@{}}
\toprule
Condition & $K$ & Intended distinction\\
\midrule
Single origin$^\star$ & 1 & root with three derivatives\\
Verbatim derivative & 1 & republication with visible source\\
Attributed paraphrase & 1 & visible source links\\
Two roots$^\star$ & 2 & two root--summary pairs\\
Shared dataset$^\star$ & 3 & two analyses reuse one dataset\\
Study and replications & 3 & summary plus two replications\\
Four independent$^\star$ & 4 & four separate samples\\
Same wording independent$^\star$ & 4 & wording without dependence\\
Same domain independent & 4 & publisher without dependence\\
Unattributed repetition & -- & repetition, provenance unknown\\
Ambiguous provenance$^\star$ & -- & unknown; abstain\\
\bottomrule
\end{tabular}
\end{table}

\paragraph{Changes from earlier releases.}
Version 1.4 addresses three problems found in our earlier releases. First, \emph{chronology}: an
audit of the v1.1 training split found \ChronReversed\ of \ChronEdges\ lineage edges whose source
was dated after its derivative (Section~\ref{sec:v11}). Generation now moves each derivative after
its sources; v1.4 has zero reversed edges. Second, \emph{observability}: a diagnostic review found
that two conditions asserted lineage that the public text did not reveal. Verbatim republications
now carry a visible upstream URL and attribution, and unattributed repetition now has an abstention
target rather than a hidden chain. Third, \emph{templating}: v1.4 has 36 topics (three times v1.1),
six prose styles per provenance role, and varied distractors. Two styles (forum and memo) appear
only in half of the test worlds, so part of the test set uses unseen surface forms.

\paragraph{Checks.}
Structural validation confirms counts, unique identifiers, verdict balance, public/gold separation,
split isolation, and chronology. A nearest-centroid probe on four coarse metadata features (mean
year, domain count, URL path shape, identifier length) predicts the 11 conditions with accuracy
\ProbeAccFour\ (chance 0.091; prespecified failure threshold 0.55). This probe cannot rule out
textual shortcuts. Labels are deterministic consequences of the generator, not human judgments.

\section{Language-Model Study}
\paragraph{Design.}
We use the six starred conditions in Table~\ref{tab:conditions} for the 12 validation worlds
(four per verdict): 72 cases per arm. No test-split world has been sent to any model. Each request contains the same four relevant documents in their public order, with title, URL,
date, outbound links, and text. Four arms differ only in their guidance:
\emph{standard} (none); \emph{reminder} (repeated reporting need not add independent evidence,
similar wording alone does not establish dependence, acknowledge unclear provenance);
\emph{EchoGraph} (reminder plus detector-proposed families, labelled as possibly wrong); and
\emph{oracle} (identical wording, gold families; for indeterminate cases, a statement that
provenance could not be determined). The oracle arm measures the value of correct metadata; it is
not a deployable method. Arm order rotates across cases.

Every arm asks for $\widehat K$. This primes all arms equally toward considering independence, so
the standard arm is a conservative baseline: it measures counting errors that remain even when the
model is explicitly asked to count.

\paragraph{Models.}
We select three model families with verified free access on the study date: Qwen3.8-27B and
GPT-OSS-120B (low reasoning effort) through Groq, and Gemini 3.5 Flash-Lite through Google's API.
Decoding uses temperature 0 and at most 2,048 output tokens. Requested and returned model
identifiers are logged. Requests rejected before generation (rate limit or unavailability) are
retried with back-off; completed responses, including invalid ones, are never re-requested.
Invalid outputs count against coverage.

\paragraph{Coverage.}
Gemini answered all 288 planned requests. The two Groq models were stopped by the free tier's
daily limit of 200,000 tokens per model (about 190 of our requests): Qwen and GPT-OSS each completed
all 24 requests for 8 of the 12 worlds, for coverage of \DevQwenCoverage\ and \DevGptossCoverage\
of planned requests in their lowest arm. Missing requests are the last worlds in plan order, not
selected by outcome. Three GPT-OSS responses cited nonexistent documents and count as invalid.

\paragraph{Analysis.}
The resampling unit is the micro-world. We report 95\% percentile intervals from 10,000
whole-world bootstrap replicates, with contrasts paired by case. With 8--12 worlds per model, the
intervals are wide and every result is exploratory: these worlds were also used to revise the
protocol (Section~\ref{sec:confirm}), so they cannot test hypotheses formed from them. We apply no
multiplicity correction and claim no confirmed effect.

\section{EchoGraph Detector}
For each time-ordered pair of documents, EchoGraph computes direct-link presence, attribution cues,
token and five-gram Jaccard similarity, bag-of-words cosine, numeric overlap, same-domain status,
and clipped time separation. A logistic classifier is trained on all pairs from determinate
training cases to predict whether the two documents share a family. Pairs above a threshold are
joined, and connected components are the families. The threshold (0.40) is chosen on validation by
lowest root-count MAE. Baselines count URLs, assign all documents to one family, group by domain,
or join pairs above a bag-of-words, TF--IDF, or MiniLM (\code{all-MiniLM-L6-v2}, pinned revision)
cosine threshold. No method uses a language model or GPU.

\section{Results}
All results in this section are from validation worlds and are exploratory.

\begin{table*}[t]
\caption{Development language-model results on validation worlds. $n$: schema-valid responses
per arm. FCR: false-corroboration rate on dependent cases (lower is better). Sens.: confidence
sensitivity, $P(\text{gold}\mid K{=}4)-P(\text{gold}\mid K{=}1)$ within world (positive is
appropriate). Ceil.: share of gold-verdict probabilities $\ge 0.99$. Under: undercount rate on
independent cases. Null$_K$: abstention on determinate cases. Abst.: abstention on ambiguous cases.
Intervals: 10,000-replicate whole-world bootstrap.}
\label{tab:llm}
\centering\footnotesize
\setlength{\tabcolsep}{3pt}
\begin{tabular}{@{}llcccccccc@{}}
\toprule
Model & Arm & $n$ & FCR [95\% CI] & Sens. [95\% CI] & Ceil. & Under & Null$_K$ & Abst. & Acc.\\
\midrule
Qwen3.8-27B & Standard & \DevQwenStandardN & \DevQwenStandardFcr\ \DevQwenStandardFcrCI & \DevQwenStandardSens\ \DevQwenStandardSensCI & \DevQwenStandardCeil & \DevQwenStandardUnder & \DevQwenStandardDetAbst & \DevQwenStandardAbst & \DevQwenStandardAcc\\
 & Reminder & \DevQwenReminderN & \DevQwenReminderFcr\ \DevQwenReminderFcrCI & \DevQwenReminderSens\ \DevQwenReminderSensCI & \DevQwenReminderCeil & \DevQwenReminderUnder & \DevQwenReminderDetAbst & \DevQwenReminderAbst & \DevQwenReminderAcc\\
 & EchoGraph & \DevQwenEchographN & \DevQwenEchographFcr\ \DevQwenEchographFcrCI & \DevQwenEchographSens\ \DevQwenEchographSensCI & \DevQwenEchographCeil & \DevQwenEchographUnder & \DevQwenEchographDetAbst & \DevQwenEchographAbst & \DevQwenEchographAcc\\
 & Oracle & \DevQwenOracleN & \DevQwenOracleFcr\ \DevQwenOracleFcrCI & \DevQwenOracleSens\ \DevQwenOracleSensCI & \DevQwenOracleCeil & \DevQwenOracleUnder & \DevQwenOracleDetAbst & \DevQwenOracleAbst & \DevQwenOracleAcc\\
\midrule
GPT-OSS-120B & Standard & \DevGptossStandardN & \DevGptossStandardFcr\ \DevGptossStandardFcrCI & \DevGptossStandardSens\ \DevGptossStandardSensCI & \DevGptossStandardCeil & \DevGptossStandardUnder & \DevGptossStandardDetAbst & \DevGptossStandardAbst & \DevGptossStandardAcc\\
 & Reminder & \DevGptossReminderN & \DevGptossReminderFcr\ \DevGptossReminderFcrCI & \DevGptossReminderSens\ \DevGptossReminderSensCI & \DevGptossReminderCeil & \DevGptossReminderUnder & \DevGptossReminderDetAbst & \DevGptossReminderAbst & \DevGptossReminderAcc\\
 & EchoGraph & \DevGptossEchographN & \DevGptossEchographFcr\ \DevGptossEchographFcrCI & \DevGptossEchographSens\ \DevGptossEchographSensCI & \DevGptossEchographCeil & \DevGptossEchographUnder & \DevGptossEchographDetAbst & \DevGptossEchographAbst & \DevGptossEchographAcc\\
 & Oracle & \DevGptossOracleN & \DevGptossOracleFcr\ \DevGptossOracleFcrCI & \DevGptossOracleSens\ \DevGptossOracleSensCI & \DevGptossOracleCeil & \DevGptossOracleUnder & \DevGptossOracleDetAbst & \DevGptossOracleAbst & \DevGptossOracleAcc\\
\midrule
Gemini 3.5 Flash-Lite & Standard & \DevGeminiStandardN & \DevGeminiStandardFcr\ \DevGeminiStandardFcrCI & \DevGeminiStandardSens\ \DevGeminiStandardSensCI & \DevGeminiStandardCeil & \DevGeminiStandardUnder & \DevGeminiStandardDetAbst & \DevGeminiStandardAbst & \DevGeminiStandardAcc\\
 & Reminder & \DevGeminiReminderN & \DevGeminiReminderFcr\ \DevGeminiReminderFcrCI & \DevGeminiReminderSens\ \DevGeminiReminderSensCI & \DevGeminiReminderCeil & \DevGeminiReminderUnder & \DevGeminiReminderDetAbst & \DevGeminiReminderAbst & \DevGeminiReminderAcc\\
 & EchoGraph & \DevGeminiEchographN & \DevGeminiEchographFcr\ \DevGeminiEchographFcrCI & \DevGeminiEchographSens\ \DevGeminiEchographSensCI & \DevGeminiEchographCeil & \DevGeminiEchographUnder & \DevGeminiEchographDetAbst & \DevGeminiEchographAbst & \DevGeminiEchographAcc\\
 & Oracle & \DevGeminiOracleN & \DevGeminiOracleFcr\ \DevGeminiOracleFcrCI & \DevGeminiOracleSens\ \DevGeminiOracleSensCI & \DevGeminiOracleCeil & \DevGeminiOracleUnder & \DevGeminiOracleDetAbst & \DevGeminiOracleAbst & \DevGeminiOracleAcc\\
\bottomrule
\end{tabular}
\end{table*}

\subsection{False corroboration (RQ1)}
Without guidance, GPT-OSS overcounts on \DevGptossStandardFcr\ of dependent cases
\DevGptossStandardFcrCI. The error has a consistent form: in 12 of its 14 overcounts the model
reported four independent sources, the number of documents. It did the same on 7 of 9 ambiguous
cases, where the correct response is to abstain. Qwen (\DevQwenStandardFcr) and Gemini
(\DevGeminiStandardFcr) rarely overcount, and both intervals include zero. Qwen nevertheless
reported four sources on 3 of 8 ambiguous cases. False corroboration therefore appears in one of
three models when provenance is visible, and it takes the form of counting documents.

\subsection{Confidence sensitivity (RQ2)}
Confidence sensitivity is indistinguishable from zero for every model and arm: a model is no more
confident in the gold verdict when four independent sources support it than when one origin does.
For Qwen and Gemini this is saturation: \DevQwenStandardCeil\ and \DevGeminiStandardCeil\ of their
gold-verdict probabilities are at least 0.99, so they could not increase. GPT-OSS stays below the
ceiling (\DevGptossStandardCeil) yet still shows no difference
(\DevGptossStandardSens\ \DevGptossStandardSensCI). In this setting, stated verdict probabilities
express how clearly the documents state the claim, not how much independent evidence supports it.

\subsection{Guidance and over-correction (RQ3)}
\begin{table*}[t]
\caption{Paired development contrasts (EchoGraph $-$ standard, and secondary contrasts). Negative
$\Delta$FCR means fewer overcounts; positive $\Delta$Under means more undercounts (worse).
Intervals: 95\% whole-world bootstrap.}
\label{tab:contrast}
\centering\small
\setlength{\tabcolsep}{8pt}
\begin{tabular}{@{}lccc@{}}
\toprule
 & Qwen & GPT-OSS & Gemini\\
\midrule
\multicolumn{4}{@{}l}{\emph{EchoGraph $-$ standard}}\\
$\Delta$FCR & \DevQwenEchoStdFcr & \DevGptossEchoStdFcr & \DevGeminiEchoStdFcr\\
95\% CI & \DevQwenEchoStdFcrCI & \DevGptossEchoStdFcrCI & \DevGeminiEchoStdFcrCI\\
$\Delta$Under & \DevQwenEchoStdUnder & \DevGptossEchoStdUnder & \DevGeminiEchoStdUnder\\
95\% CI & \DevQwenEchoStdUnderCI & \DevGptossEchoStdUnderCI & \DevGeminiEchoStdUnderCI\\
$\Delta$Acc & \DevQwenEchoStdAcc & \DevGptossEchoStdAcc & \DevGeminiEchoStdAcc\\
95\% CI & \DevQwenEchoStdAccCI & \DevGptossEchoStdAccCI & \DevGeminiEchoStdAccCI\\
\midrule
\multicolumn{4}{@{}l}{\emph{Secondary $\Delta$FCR}}\\
Reminder $-$ std. & \DevQwenRemStdFcr & \DevGptossRemStdFcr & \DevGeminiRemStdFcr\\
EchoGraph $-$ rem. & \DevQwenEchoRemFcr & \DevGptossEchoRemFcr & \DevGeminiEchoRemFcr\\
Oracle $-$ EchoGraph & \DevQwenOracleEchoFcr & \DevGptossOracleEchoFcr & \DevGeminiOracleEchoFcr\\
\midrule
Coverage (min.\ arm) & \DevQwenCoverage & \DevGptossCoverage & \DevGeminiCoverage\\
\bottomrule
\end{tabular}
\end{table*}

Every form of guidance eliminates GPT-OSS's overcounting (Table~\ref{tab:contrast}), but the plain
reminder does so as completely as EchoGraph families do (EchoGraph $-$ reminder:
\DevGptossEchoRemFcr). The reduction comes with a large increase in undercounting
(\DevGptossEchoStdUnder\ \DevGptossEchoStdUnderCI): with guidance, GPT-OSS reports a single source
for most cases, including independent ones. This persists in the oracle arm, where the prompt lists
the four correct singleton families: on 13 of 18 independent cases GPT-OSS still reported one
source (undercount \DevGptossOracleUnder). The model does not appear to use the metadata; instead,
the instruction to consider dependence seems to shift its default answer from ``count the
documents'' to ``one origin''. Gemini over-corrects differently. The counts it gives remain
accurate, but with guidance it abstains on determinate cases (\DevGeminiEchographDetAbst\ with
EchoGraph versus \DevGeminiStandardDetAbst\ without), mostly on independent ones: 9 of its 13
abstentions in the EchoGraph arm were on four-independent or same-wording cases. The reminder also lowered Gemini's verdict accuracy
(\DevGeminiRemStdAcc). Qwen counted almost perfectly in every arm, leaving nothing for guidance to
fix.

Our protocol's original mitigation criterion checked only verdict accuracy. By that criterion,
GPT-OSS would have shown a large, apparently successful reduction in false corroboration. A
mitigation for overcounting must therefore be evaluated together with undercounting and abstention
on independent evidence.

\subsection{Detector results}
\begin{table}[t]
\caption{Source-family detection on v1.4 determinate cases. Validation selected the thresholds;
test is a single run after the protocol freeze (648 cases, 72 worlds). Higher F1 and lower MAE are
better.}
\label{tab:detection}
\centering\small
\begin{tabular}{@{}lcccc@{}}
\toprule
 & \multicolumn{2}{c}{Validation} & \multicolumn{2}{c}{Test}\\
Method & F1 & MAE & F1 & MAE\\
\midrule
URL singletons & \UrlFOneVal & \UrlMAEVal & \UrlFOneTest & \UrlMAETest\\
All one family & \OneFOneVal & \OneMAEVal & \OneFOneTest & \OneMAETest\\
Domain & \DomainFOneVal & \DomainMAEVal & \DomainFOneTest & \DomainMAETest\\
Bag of words & \BowFOneVal & \BowMAEVal & \BowFOneTest & \BowMAETest\\
TF--IDF & \TfidfFOneVal & \TfidfMAEVal & \TfidfFOneTest & \TfidfMAETest\\
MiniLM & \MinilmFOneVal & \MinilmMAEVal & \MinilmFOneTest & \MinilmMAETest\\
\textbf{EchoGraph} & \textbf{\EchoFOneVal} & \textbf{\EchoMAEVal} & \textbf{\EchoFOneTest} & \textbf{\EchoMAETest}\\
\quad no links & \NoLinksFOneVal & \NoLinksMAEVal & \NoLinksFOneTest & \NoLinksMAETest\\
\bottomrule
\end{tabular}
\end{table}

On validation, EchoGraph reaches pair F1 \EchoFOneVal\ and MAE \EchoMAEVal, against at most
\OneFOneVal\ F1 for any baseline (Table~\ref{tab:detection}). Removing link features lowers F1 to
\NoLinksFOneVal: once provenance is made observable, much of the task is reading explicit links and
attribution. Shared datasets remain the main failure (F1 \EchoSharedFOneVal), because a reused
dataset is signalled by an identifier rather than by a link between documents. The single held-out
run after the protocol freeze agrees with validation: EchoGraph reaches F1 \EchoFOneTest\
\EchoFOneTestCI\ and MAE \EchoMAETest, and shared datasets remain the failure (F1
\EchoSharedFOneTest). High detector accuracy is a property of this construction, not evidence of
open-web provenance recovery.

\subsection{From development to a confirmatory test}\label{sec:confirm}
These observations changed the held-out protocol, which was then frozen by SHA-256 before any
test-split access. Mitigation now
requires, in addition to lower FCR and verdict accuracy within 0.05 of the standard arm, that
undercounting on independent cases rise by no more than 0.10; models whose standard-arm FCR is zero
are reported as untestable rather than as failures; and confidence sensitivity is reported with
the ceiling rate. The amended protocol covers all 72 test worlds (1,728 requests per model). Under
the free daily token limit, it requires about ten days per Groq model, which is why this paper
reports development results only.

\subsection{What the v1.1 chronology defect taught us}\label{sec:v11}
Our earlier release (v1.1; 2,640 cases, 240 worlds) was evaluated once under a frozen protocol.
EchoGraph obtained held-out pair F1 \EchoFOne\ \EchoFOneCI\ and MAE \EchoMAE, against \UrlFOne\ for
URL counting. A training-only audit after the freeze then found that \ChronReversed\ of
\ChronEdges\ lineage edges (26\%) ran backwards in time. Because EchoGraph orders document pairs by
date before computing link features, the defect could affect the detector itself, not only its
chronology feature, so those scores cannot be carried over to corrected data. We preserved the
frozen v1.1 artifacts unchanged rather than silently regenerating them, and the defect motivated
the repairs in v1.4. The episode illustrates that structural validators for synthetic provenance
benchmarks should check temporal consistency of lineage, not only counts and splits.

\section{Reproducibility}
Code, data, and logs are available at
\url{https://github.com/theaayushdev/echotrace-benchmark}. The repository contains the deterministic generators, public/gold data separation, validators,
detector code and validation-selected settings, the frozen confirmatory protocol,
the study runner, raw append-only response logs including failures, and analysis scripts. Every
table value in this paper is generated from analysis artifacts by
\code{scripts/export\_llm\_study\_results.py} and \code{scripts/export\_v1\_paper\_results.py};
response counts quoted in the text come from \code{scripts/llm\_count\_breakdown.py}. The study requires no GPU and no paid API; it runs on
CPU and free provider tiers. Raw logs contain no credentials.

\section{Conclusion}
In closed packets with visible provenance, one of three freely accessible models counted documents
as independent sources, and no model's verdict confidence depended on source independence. Asking
models to consider dependence removed overcounting but caused undercounting or abstention on
genuinely independent evidence, even with correct provenance metadata. Provenance mitigations
should therefore be judged on both directions of counting error. The next steps are the held-out
run of the frozen protocol, blinded human review of the free-text answers, and real-source
evaluation with retrieval.

\section*{Limitations}\label{sec:limits}
\paragraph{Synthetic construct.}
Documents are generated from templates and labels follow generator rules. In v1.4 provenance is
usually visible through links and attribution, so the study measures whether models \emph{use}
visible provenance, not whether they discover hidden ancestry. Results on unseen prose styles
reduce, but do not remove, the risk of template learning.

\paragraph{Exploratory evidence.}
All language-model results come from 12 validation worlds (8 complete worlds for the Groq models),
which also informed the protocol revision. Intervals are wide, no multiplicity correction is
applied, and no result should be read as a confirmed effect.

\paragraph{No human validation.}
No independent human annotation is complete. Primary outcomes are computed from a structured count,
not from human judgments of whether free-text answers claim corroboration. The answers are retained
for a later blinded human rubric.

\paragraph{Closed evidence packets.}
Models receive the four relevant documents. We do not evaluate retrieval, distractor rejection, or
an autonomous agent. Connected components force non-overlapping families, whereas real articles may
draw on several studies.

\paragraph{Models and access.}
We use three openly accessible models under free quotas. Provider aliases may change after the
study date, and results may not transfer to larger proprietary systems. Asking for $\widehat K$ may
itself change behaviour; we treat it as part of the task definition, applied equally to all arms.

\paragraph{Responsible use.}
Dependence is not a proxy for truth, quality, or misconduct. Similar wording does not establish
copying, and a shared family can be reliable. EchoTrace should not be used to accuse authors or rank
publishers.

\bibliography{references}
\end{document}
